\documentclass[aspectratio=169]{beamer}

\usepackage[utf8]{inputenc}
\usepackage[T1]{fontenc}
\usepackage{amsmath,amssymb,amsfonts,bm,mathtools}
\usepackage{physics}
\usepackage{braket}
\usepackage{quantikz}
\usepackage{hyperref}
\usepackage{booktabs}

\usetheme{Madrid}

\title{Quantum Approximate Optimization Algorithm (QAOA)}
\subtitle{From Combinatorial Optimization to Quantum Neural Networks}
\author{Morten Hjorth-Jensen}
\date{Spring 2026}

\begin{document}

\frame{\titlepage}

% ─────────────────────────────────────────────────────────────────────────────
\begin{frame}{Outline}
\tableofcontents
\end{frame}

% =============================================================================
\section{Motivation}
% =============================================================================

\begin{frame}{Why QAOA?}
Many optimization problems can be written as the minimization or maximization of a
cost function over binary variables,
\[
x_i \in \{0,1\}, \qquad i=1,\dots,n.
\]

Examples:
\begin{itemize}
    \item MaxCut,
    \item portfolio optimization,
    \item scheduling,
    \item constraint satisfaction,
    \item graph partitioning.
\end{itemize}

QAOA is a variational quantum algorithm designed to attack such combinatorial
optimization problems by encoding the cost function into a quantum Hamiltonian.
\end{frame}

\begin{frame}{High-Level Idea}
QAOA constructs a parameterized quantum state
\[
\ket{\psi_p(\bm{\gamma},\bm{\beta})}
=
\prod_{\ell=1}^p e^{-i\beta_\ell H_M} e^{-i\gamma_\ell H_C}\ket{+}^{\otimes n},
\]
where
\begin{itemize}
    \item \(H_C\) is the cost Hamiltonian (encodes the objective),
    \item \(H_M\) is the mixer Hamiltonian (generates transitions),
    \item \(p\) is the circuit depth,
    \item \(\bm{\gamma}=(\gamma_1,\dots,\gamma_p)\), \(\bm{\beta}=(\beta_1,\dots,\beta_p)\).
\end{itemize}

The parameters are optimized classically to maximize \(\ev{H_C}{\psi_p}\).
\end{frame}

\begin{frame}{Goals of These Notes}
We develop four themes:
\begin{enumerate}
    \item rigorous derivation of the QAOA cost Hamiltonian from combinatorial
          optimization via QUBO and Ising models,
    \item analytical results for \(p=1\) QAOA applied to MaxCut,
    \item the interpretation of QAOA as a discretized adiabatic protocol,
    \item explicit connections between QAOA and quantum neural networks (QNNs).
\end{enumerate}

\medskip
The central idea is that QAOA is simultaneously:
\begin{itemize}
    \item an optimization algorithm for combinatorial problems,
    \item a discretized adiabatic protocol,
    \item and a structured parameterized quantum circuit / QNN.
\end{itemize}
\end{frame}

% =============================================================================
\section{QUBO and the Ising Mapping}
% =============================================================================

\begin{frame}{Rigorous Formulation of QUBO}
A quadratic unconstrained binary optimization (QUBO) problem is
\[
\min_{x\in\{0,1\}^n} f(x),
\]
with
\[
f(x)=x^T Q x + c^T x + c_0,
\]
where
\begin{itemize}
    \item \(x=(x_1,\dots,x_n)^T\),
    \item \(Q\in\mathbb{R}^{n\times n}\),
    \item \(c\in\mathbb{R}^n\),
    \item \(c_0\in\mathbb{R}\).
\end{itemize}

Because \(x_i^2=x_i\) for binary variables, any quadratic binary polynomial
can be reduced to this form.
\end{frame}

\begin{frame}{Expanded QUBO Form}
Writing out the quadratic form,
\[
f(x)=\sum_{i,j=1}^n Q_{ij}x_i x_j + \sum_{i=1}^n c_i x_i + c_0.
\]

If \(Q\) is symmetrized, one may equivalently write
\[
f(x)=\sum_{i=1}^n a_i x_i + \sum_{1\le i<j\le n} b_{ij} x_i x_j + c_0.
\]

Constraint terms are often converted into QUBO penalties, for example
\[
\lambda\Big(\sum_i x_i - k\Big)^2,
\]
which enforces \(\sum_i x_i \approx k\) when \(\lambda\) is sufficiently large.
\end{frame}

\begin{frame}{Why QUBO is So Central}
QUBO is important because many discrete optimization problems can be transformed
into it.

\[
\text{optimization problem}
\;\longrightarrow\;
\text{binary variables}
\;\longrightarrow\;
\text{quadratic objective}
\;\longrightarrow\;
\text{Ising Hamiltonian}
\;\longrightarrow\;
\text{quantum cost Hamiltonian}.
\]

QAOA acts on the final object in this chain.
\end{frame}

\begin{frame}{Binary-to-Spin Transformation}
To map QUBO to an Ising model, introduce spin variables
\[
z_i \in \{-1,+1\}
\]
through
\[
x_i = \frac{1-z_i}{2},
\qquad
z_i = 1-2x_i.
\]

\textbf{Linear terms:}
\[
a_i x_i = \frac{a_i}{2} - \frac{a_i}{2} z_i.
\]

\textbf{Quadratic terms:}
\[
b_{ij} x_i x_j
=
\frac{b_{ij}}{4}\bigl(1 - z_i - z_j + z_i z_j\bigr).
\]

Every quadratic binary interaction becomes a constant, local fields, and a
pairwise Ising coupling.
\end{frame}

\begin{frame}{General Ising Form}
After the substitution \(x_i=(1-z_i)/2\), any QUBO becomes
\[
f(x)
=
\mathrm{const}
+
\sum_i h_i z_i
+
\sum_{i<j} J_{ij} z_i z_j.
\]

The quantum cost Hamiltonian is then obtained by promoting classical spins to
Pauli-\(Z\) operators,
\[
z_i \;\longrightarrow\; Z_i,
\qquad
z_i z_j \;\longrightarrow\; Z_i Z_j,
\]
giving
\[
H_C
=
\sum_i h_i Z_i
+
\sum_{i<j} J_{ij} Z_i Z_j
+
\mathrm{const}.
\]

This Hamiltonian is diagonal in the computational basis.
\end{frame}

% =============================================================================
\section{MaxCut: From Graph to Quantum Hamiltonian}
% =============================================================================

\begin{frame}{Definition of MaxCut}
Let \(G=(V,E)\) be a graph with edge weights \(w_{ij}\ge 0\).
A cut partitions the vertices into two sets.

For binary variables \(x_i\in\{0,1\}\), interpret
\[
x_i =
\begin{cases}
0, & i \in S,\\
1, & i \in \bar S.
\end{cases}
\]

An edge \((i,j)\) is cut if \(x_i\neq x_j\). The weighted MaxCut objective is
\[
\max_{x\in\{0,1\}^n}
C(x)
=
\sum_{(i,j)\in E} w_{ij}\,\mathbf{1}_{x_i\neq x_j}.
\]
\end{frame}

\begin{frame}{MaxCut as a QUBO and Ising Hamiltonian}
Because \(\mathbf{1}_{x_i\neq x_j} = x_i+x_j-2x_ix_j\), the cut objective is
\[
C(x)
=
\sum_{(i,j)\in E} w_{ij}\bigl[x_i+x_j-2x_i x_j\bigr].
\]

After the spin substitution \(x_i=(1-z_i)/2\), one finds
\[
x_i+x_j-2x_ix_j = \frac{1}{2}(1-z_iz_j).
\]

Therefore,
\[
C(z)
=
\sum_{(i,j)\in E} \frac{w_{ij}}{2}(1-z_iz_j).
\]

Promoting \(z_i\) to Pauli operators gives the MaxCut cost Hamiltonian:
\[
H_C
=
\sum_{(i,j)\in E} \frac{w_{ij}}{2}(I-Z_iZ_j).
\]
\end{frame}

\begin{frame}{Why This Hamiltonian Encodes MaxCut}
On a computational basis state \(\ket{z_1\cdots z_n}\), each \(Z_i\) has
eigenvalue \(z_i\in\{-1,+1\}\). Hence
\[
Z_i Z_j \ket{z} = z_i z_j \ket{z}.
\]

Therefore
\[
\frac{1-z_i z_j}{2}
=
\begin{cases}
0, & z_i=z_j,\\
1, & z_i\neq z_j.
\end{cases}
\]

So each edge contributes exactly its cut value \(w_{ij}\frac{1-z_iz_j}{2}\).
\end{frame}

\begin{frame}{Worked MaxCut Example: Triangle Graph}
Consider a 3-vertex complete graph with unit weights:
\[
E=\{(1,2),(2,3),(1,3)\}.
\]

Then
\[
H_C
=
\frac{1}{2}(I-Z_1Z_2)
+
\frac{1}{2}(I-Z_2Z_3)
+
\frac{1}{2}(I-Z_1Z_3)
=
\frac{3}{2}I - \frac{1}{2}(Z_1Z_2+Z_2Z_3+Z_1Z_3).
\]

The maximum cut has value \(2\), achieved by any state with one spin
separated from the other two.
\end{frame}

% =============================================================================
\section{QAOA Hamiltonians and the QAOA State}
% =============================================================================

\begin{frame}{Cost Hamiltonian}
For a generic problem, define \(H_C\) so that computational basis states are
eigenstates with the classical objective as eigenvalues:
\[
H_C\ket{x}=C(x)\ket{x}.
\]

For MaxCut,
\[
H_C = \sum_{(i,j)\in E} w_{ij}\frac{I-Z_iZ_j}{2}.
\]

The QAOA algorithm maximizes
\[
\expval{H_C}.
\]
\end{frame}

\begin{frame}{Mixer Hamiltonian}
The standard choice of mixer is
\[
H_M = \sum_{i=1}^n X_i.
\]

Starting from
\[
\ket{+}^{\otimes n}
=
\bigotimes_{i=1}^n \frac{\ket{0}+\ket{1}}{\sqrt{2}},
\]
the mixer drives transitions between computational basis states:
\[
X_i\ket{x_1\cdots x_i\cdots x_n}
=
\ket{x_1\cdots (1-x_i)\cdots x_n}.
\]

QAOA requires
\[
[H_C,H_M]\neq 0
\]
to generate a nontrivial variational manifold.
\end{frame}

\begin{frame}{The QAOA State}
At depth \(p\), the QAOA state is
\[
\ket{\psi_p(\bm{\gamma},\bm{\beta})}
=
e^{-i\beta_p H_M}e^{-i\gamma_p H_C}
\cdots
e^{-i\beta_1 H_M}e^{-i\gamma_1 H_C}
\ket{+}^{\otimes n}.
\]

The variational objective is
\[
F_p(\bm{\gamma},\bm{\beta})
=
\ev{H_C}{\psi_p(\bm{\gamma},\bm{\beta})}.
\]

Parameters \((\bm{\gamma}^\star,\bm{\beta}^\star)\) are found by the classical
outer loop:
\[
(\bm{\gamma}^\star,\bm{\beta}^\star)
=
\arg\max_{\bm{\gamma},\bm{\beta}}
F_p(\bm{\gamma},\bm{\beta}).
\]
\end{frame}

% =============================================================================
\section{From Adiabatic Evolution to QAOA}
% =============================================================================

\begin{frame}{Adiabatic Optimization and Its Discretization}
Start from the adiabatic interpolation
\[
H(s) = (1-s)H_M + s H_C,
\qquad s\in[0,1].
\]

The ground state of \(H(0)=H_M\) is \(\ket{+}^{\otimes n}\) (easy to prepare).
The ground state of \(H(1)=H_C\) encodes the optimal solution.

Discretizing into \(p\) time slices and applying first-order Trotterization,
\[
e^{-i\Delta t_\ell\left[(1-s_\ell)H_M+s_\ell H_C\right]}
\approx
e^{-i\beta_\ell H_M}e^{-i\gamma_\ell H_C},
\]
with \(\beta_\ell=(1-s_\ell)\Delta t_\ell\), \(\gamma_\ell=s_\ell\Delta t_\ell\), yields:
\[
\ket{\psi_p(\bm{\gamma},\bm{\beta})}
=
\prod_{\ell=1}^p e^{-i\beta_\ell H_M}e^{-i\gamma_\ell H_C}
\ket{+}^{\otimes n}.
\]
\end{frame}

\begin{frame}{Why the Parameters Become Variational}
In strict adiabatic evolution, the schedule \(s(t)\) is fixed beforehand.

In QAOA, the effective step parameters
\[
\{(\gamma_\ell,\beta_\ell)\}_{\ell=1}^p
\]
are treated as \emph{free variational parameters} and optimized classically.

This changes the problem from
\[
\text{simulate a prescribed schedule}
\]
to
\[
\text{optimize over a family of Trotterized schedules}.
\]

QAOA is therefore a discretized, variationally optimized form of adiabatic
evolution.
\end{frame}

\begin{frame}{Variational Manifold}
Define the QAOA manifold
\[
\mathcal{M}_p
=
\left\{
\ket{\psi_p(\bm{\gamma},\bm{\beta})}
:
\bm{\gamma},\bm{\beta}\in\mathbb{R}^p
\right\}.
\]

QAOA searches only within \(\mathcal{M}_p\). Its expressive power depends on:
\begin{itemize}
    \item the depth \(p\),
    \item the choice of cost Hamiltonian,
    \item the choice of mixer,
    \item the optimizer used in the classical loop.
\end{itemize}

As \(p\to\infty\) with suitable schedules, \(\mathcal{M}_p\) can approximate the
adiabatic solution arbitrarily well.
\end{frame}

% =============================================================================
\section{Circuit-Level Implementation}
% =============================================================================

\begin{frame}{Cost Unitary for MaxCut}
Because all \(Z_iZ_j\) terms commute, the cost unitary factorizes:
\[
e^{-i\gamma H_C}
=
\prod_{(i,j)\in E}
e^{-i\gamma w_{ij}(I-Z_iZ_j)/2}.
\]

Each edge corresponds to a two-qubit \(ZZ\)-phase interaction, implementable as
a standard \texttt{CNOT}--\(R_z\)--\texttt{CNOT} sequence.

The mixer unitary also factorizes over qubits since \([X_i,X_j]=0\):
\[
e^{-i\beta H_M}
=
\prod_{i=1}^n e^{-i\beta X_i}
=
\prod_{i=1}^n R_x(2\beta).
\]
\end{frame}

\begin{frame}{Quantikz: One MaxCut Edge}
For one edge \((i,j)\) with weight \(w\), the cost unitary
\(e^{-i\gamma \frac{w}{2}(I-Z_iZ_j)}\)
is implemented (up to global phase) by a \(ZZ\)-rotation:

\begin{quantikz}[row sep=0.3cm]
\lstick{$\ket{q_i}$} & \ctrl{1} & \qw                   & \ctrl{1} & \qw \\
\lstick{$\ket{q_j}$} & \targ{}  & \gate{R_z(\gamma w)}  & \targ{}  & \qw
\end{quantikz}
\end{frame}

\begin{frame}{Quantikz: Mixer Layer}
Because \(e^{-i\beta X} = R_x(2\beta)\), the full mixer layer
\(e^{-i\beta H_M}=\prod_i e^{-i\beta X_i}\)
is simply a product of single-qubit \(X\)-rotations:

\begin{quantikz}[row sep=0.25cm]
\lstick{$\ket{q_1}$} & \gate{R_x(2\beta)} & \qw \\
\lstick{$\ket{q_2}$} & \gate{R_x(2\beta)} & \qw \\
\lstick{$\ket{q_3}$} & \gate{R_x(2\beta)} & \qw \\
\lstick{$\vdots$}    & \vdots             &     \\
\lstick{$\ket{q_n}$} & \gate{R_x(2\beta)} & \qw
\end{quantikz}
\end{frame}

\begin{frame}{Quantikz: One Full \(p=1\) QAOA Layer}
For a 3-qubit line graph with edges \((1,2)\) and \((2,3)\), one layer looks like:

\begin{quantikz}[row sep=0.28cm, column sep=0.38cm]
\lstick{$\ket{+}$} & \ctrl{1} & \qw & \ctrl{1} & \gate{R_x(2\beta)} & \qw \\
\lstick{$\ket{+}$} & \targ{}  & \gate{R_z(\gamma w_{12})} & \targ{} & \ctrl{1} & \gate{R_x(2\beta)} \\
\lstick{$\ket{+}$} & \qw      & \qw & \qw      & \targ{} & \gate{R_z(\gamma w_{23})}
\end{quantikz}

\medskip
The layer alternates:
\[
\underbrace{\text{cost layer}}_{\text{problem-dependent entangling phases}}
\;\longrightarrow\;
\underbrace{\text{mixer layer}}_{\text{global single-qubit rotations}}.
\]
\end{frame}

% =============================================================================
\section{Analytical \texorpdfstring{\(p=1\)}{p=1} Results for MaxCut}
% =============================================================================

\begin{frame}{The \(p=1\) Objective}
At depth \(p=1\),
\[
\ket{\psi_1(\gamma,\beta)}
=
e^{-i\beta H_M}e^{-i\gamma H_C}\ket{+}^{\otimes n},
\]
and the optimization problem is
\[
\max_{\gamma,\beta}
F_1(\gamma,\beta)
=
\ev{H_C}{\psi_1(\gamma,\beta)}.
\]

Because \(H_C=\sum_{(i,j)\in E}C_{ij}\), this reduces to a sum over edges:
\[
F_1(\gamma,\beta)
=
\sum_{(i,j)\in E} \ev{C_{ij}}{\psi_1}.
\]

The key object to compute is therefore \(\ev{Z_iZ_j}_{\gamma,\beta}\).
\end{frame}

\begin{frame}{Heisenberg-Picture Strategy}
Instead of evolving the state, evolve the operator:
\[
\ev{Z_iZ_j}_{\gamma,\beta}
=
\bra{+}^{\otimes n}
e^{i\gamma H_C}
e^{i\beta H_M}
Z_iZ_j
e^{-i\beta H_M}
e^{-i\gamma H_C}
\ket{+}^{\otimes n}.
\]

Because \(\ket{+}\) is an eigenstate of \(X\), only terms that are products
of \(X\)-operators survive:
\[
\bra{+} X \ket{+}=1,\qquad
\bra{+} Y \ket{+}=0,\qquad
\bra{+} Z \ket{+}=0.
\]

\textbf{Conjugation by the mixer:}
\[
e^{i\beta H_M} Z_i e^{-i\beta H_M}
=
Z_i\cos(2\beta)+Y_i\sin(2\beta).
\]

So \(e^{i\beta H_M}Z_iZ_je^{-i\beta H_M}\) expands into four terms involving
\(Z_iZ_j\), \(Z_iY_j\), \(Y_iZ_j\), and \(Y_iY_j\).
\end{frame}

\begin{frame}{Conjugation by the Cost Hamiltonian}
Because \(H_C\) contains only \(Z\)-type terms, conjugation by
\(e^{i\gamma H_C}\) can be analyzed locally.

For a single neighbor interaction,
\[
e^{i\theta Z_iZ_j} Y_i e^{-i\theta Z_iZ_j}
=
Y_i\cos(2\theta)+X_i Z_j\sin(2\theta),
\]
and similarly for \(Y_j\).

These identities convert \(Y\) operators into \(X\) operators, which can survive
in the \(\ket{+}^{\otimes n}\) expectation value.

For \(Y_iY_j\), one finds
\[
e^{i\theta Z_iZ_j} Y_iY_j e^{-i\theta Z_iZ_j}
= Y_iY_j\cos^2(2\theta) - X_iX_j\sin^2(2\theta) + \cdots.
\]
\end{frame}

\begin{frame}{Closed-Form Result: Unweighted \(d\)-Regular Graphs}
For an unweighted \(d\)-regular, triangle-free graph, all edges are equivalent
and the \(p=1\) expectation value has the closed form:
\[
\ev{C_{ij}}
=
\frac{1}{2}
+
\frac{1}{2}\sin(4\beta)\sin\gamma \cos^{d-1}\gamma.
\]

Hence the total expected cut value is
\[
\boxed{
F_1(\gamma,\beta)
=
|E|
\left[
\frac{1}{2}
+
\frac{1}{2}\sin(4\beta)\sin\gamma \cos^{d-1}\gamma
\right].
}
\]

\begin{itemize}
    \item \(\sin(4\beta)\) measures mixer-generated interference efficiency.
    \item \(\sin\gamma\) generates the edge-dependent cost phase.
    \item \(\cos^{d-1}\gamma\) encodes the suppression from \(d-1\) neighboring edges.
\end{itemize}
\end{frame}

\begin{frame}{Optimizing the \(p=1\) Formula}
The formula is maximized by choosing \(\sin(4\beta)=1\), i.e.\ \(\beta=\pi/8\),
which reduces the problem to maximizing
\[
f(\gamma)=\sin\gamma \cos^{d-1}\gamma.
\]

Differentiating,
\[
f'(\gamma)=\cos^d\gamma - (d-1)\sin^2\gamma\,\cos^{d-2}\gamma = 0,
\]
gives
\[
\tan^2\gamma = \frac{1}{d-1},
\qquad
\tan\gamma = \frac{1}{\sqrt{d-1}}.
\]

For \(d=3\) this yields the well-known approximation ratio
\[
\alpha_{p=1} \approx 0.6924,
\]
already nontrivially above the random-cut bound of \(1/2\).
\end{frame}

\begin{frame}{Example: Ring Graph (\(d=2\))}
For a cycle graph (\(d=2\)):
\[
F_1(\gamma,\beta)
=
|E|
\left[
\frac{1}{2}
+
\frac{1}{4}\sin(4\beta)\sin(2\gamma)
\right].
\]

The optimum is achieved for
\[
\beta=\frac{\pi}{8},
\qquad
\gamma=\frac{\pi}{4},
\]
giving
\[
F_1^{\max} = \frac{3|E|}{4}.
\]

So QAOA at \(p=1\) recovers \(75\%\) of the maximum cut on a ring,
using just two parameters.
\end{frame}

\begin{frame}{Expressivity vs Depth}
Increasing \(p\) enlarges the variational manifold:
\[
\mathcal{M}_1 \;\subset\; \mathcal{M}_2 \;\subset\; \cdots
\]

In the limit \(p\to\infty\) one expects convergence to the adiabatic solution
for suitable schedules.

But larger \(p\) also means:
\begin{itemize}
    \item more parameters to optimize,
    \item harder classical optimization (barren plateaus),
    \item deeper circuits,
    \item more hardware noise.
\end{itemize}

The \(p=1\) results above show that even the shallowest layer captures
nontrivial graph structure.
\end{frame}

% =============================================================================
\section{QAOA and Quantum Neural Networks}
% =============================================================================

\begin{frame}{QAOA as a Parameterized Quantum Circuit}
A general quantum neural network (QNN) or parameterized quantum circuit has
the form
\[
\ket{\psi(\bm{\theta})}=U(\bm{\theta})\ket{\psi_0}.
\]

QAOA is a highly structured special case:
\[
U(\bm{\theta})
=
\prod_{\ell=1}^p e^{-i\beta_\ell H_M}e^{-i\gamma_\ell H_C},
\qquad
\bm{\theta}=(\bm{\gamma},\bm{\beta}).
\]

In a classical neural network one has
\(x\mapsto W_L\sigma(\cdots\sigma(W_1 x))\);
in a QNN one has alternating parameterized unitaries.
QAOA fits exactly this architecture, where each QAOA layer is a structured
quantum neural layer.
\end{frame}

\begin{frame}{QAOA, VQE, and QNNs}
These three frameworks are closely related:

\medskip
\begin{center}
\begin{tabular}{ll}
\toprule
\textbf{VQE:}  & \(\displaystyle\min_{\bm{\theta}} \ev{H}{\psi(\bm{\theta})}\) \\[8pt]
\textbf{QAOA:} & \(\displaystyle\max_{\bm{\gamma},\bm{\beta}} \ev{H_C}{\psi_p(\bm{\gamma},\bm{\beta})}\) \\[8pt]
\textbf{QNN:}  & \(\displaystyle\min_{\bm{\theta}} \mathcal{L}\bigl(\ev{O_1},\ev{O_2},\dots\bigr)\) \\
\bottomrule
\end{tabular}
\end{center}

\medskip
All are instances of:
\begin{itemize}
    \item a parameterized quantum circuit,
    \item a classical loss or objective function,
    \item a hybrid quantum-classical optimization loop.
\end{itemize}
\end{frame}

\begin{frame}{Task-Informed vs Generic QNNs}
QAOA differs from a generic hardware-efficient ansatz because its architecture
is not arbitrary.

Instead, it is designed around:
\begin{itemize}
    \item the problem Hamiltonian \(H_C\) (problem-specific entangling structure),
    \item a physically motivated mixer \(H_M\),
    \item an optimization objective directly aligned with the task.
\end{itemize}

So QAOA is a \textbf{task-informed QNN}, not a black-box trainable circuit.

\medskip
Gradients can be computed via the \textbf{parameter-shift rule}:
\[
\frac{\partial}{\partial\gamma_\ell}F_p
=
\frac{F_p(\gamma_\ell+\pi/2)-F_p(\gamma_\ell-\pi/2)}{2},
\]
enabling gradient-based classical optimization without finite differences.
\end{frame}

\begin{frame}{Variational Principle and Expressivity}
All three frameworks share the same variational structure:
\begin{enumerate}
    \item a \textbf{manifold of trial states},
    \item an \textbf{objective functional}.
\end{enumerate}

\medskip
\begin{center}
\begin{tabular}{lll}
\toprule
Method & Manifold & Objective \\
\midrule
QAOA & \(\mathcal{M}_p = \{\ket{\psi_p(\bm{\gamma},\bm{\beta})}\}\) & \(\ev{H_C}{\psi_p}\) \\[4pt]
VQE  & \(\{\ket{\psi(\bm{\theta})}\}\) & \(\ev{H}{\psi}\) \\[4pt]
QNN  & \(\{\ket{\psi(\bm{\theta})}\}\) & task-dependent loss \\
\bottomrule
\end{tabular}
\end{center}

\medskip
QAOA is best viewed as a \textbf{variational principle with a structured and
interpretable ansatz} shaped by the problem's combinatorial structure.
\end{frame}

% =============================================================================
\section{Summary}
% =============================================================================

\begin{frame}{Conceptual Summary: The Two Derivation Chains}
QAOA sits at the intersection of two independent derivation chains:

\medskip
\textbf{From optimization:}
\[
\text{combinatorial opt.}
\;\to\;
\text{QUBO}
\;\to\;
\text{Ising model}
\;\to\;
\text{Pauli Hamiltonian}
\;\to\;
H_C.
\]

\textbf{From adiabatic physics:}
\[
\text{adiabatic evolution}
\;\to\;
\text{Trotterization}
\;\to\;
\text{variational schedule}
\;\to\;
\text{QAOA ansatz}.
\]

\medskip
Both chains arrive at the same object:
\[
\ket{\psi_p(\bm{\gamma},\bm{\beta})}
=
\prod_{\ell=1}^p e^{-i\beta_\ell H_M}e^{-i\gamma_\ell H_C}\ket{+}^{\otimes n}.
\]
\end{frame}

\begin{frame}{Final Summary}
We have developed:
\begin{itemize}
    \item a rigorous formulation of QUBO and its mapping to Ising / Pauli Hamiltonians,
    \item the MaxCut Hamiltonian \(H_C=\sum_{(i,j)\in E}\frac{w_{ij}}{2}(I-Z_iZ_j)\)
          derived step by step,
    \item explicit definitions of the QAOA cost and mixer Hamiltonians,
    \item a derivation of QAOA from adiabatic evolution via Trotterization,
    \item closed-form \(p=1\) analytical results: approximation ratio
          \(\alpha_{p=1}\approx 0.6924\) for 3-regular graphs,
    \item circuit-level realizations using quantikz,
    \item the connection between QAOA and quantum neural networks / VQE.
\end{itemize}

\medskip
QAOA is therefore a structured variational quantum algorithm for discrete
optimization with deep links to physics and quantum machine learning.
\end{frame}

\end{document}
